\section{Related Work}
\label{sec:related_work}
\subsection{Radiance Fields}
Neural Radiance Field (NeRF) \cite{nerf} has arisen as a significant development in the field of Computer Vision and Computer Graphics, used for synthesizing novel views of a scene from a sparse set of images by combining machine learning with geometric reasoning. Recently, a plethora of research and methodologies built upon NeRF have been proposed. For example, \cite{park2021nerfies,pumarola2021dnerf,li2021SceneFlownerf,yan2023nerfds} extend NeRF to dynamic and non-grid scenes, \cite{barron2021mipnerf,barron2022mipnerf360} significantly improve the rendering quality of NeRF. Recently, researchers have collectively recognized that the bottleneck in efficiency lies in querying the neural field, prompting efforts to address it. InstantNGP \cite{muller2022instant} combines a neural network with a multiresolution
hash table for efficient evaluation while Plenoxels \cite{fridovich2022plenoxels} replaces neural networks with a sparse voxel grid. 3D Gaussian splatting \cite{3dgs} further adopts a discrete 3D Gaussian representation of scenes, significantly accelerating the training and rendering of radiance fields. It has attracted considerable
research interest in the field of generation \cite{yi2023gaussiandreamer,chen2023gaussianeditor}, relighting \cite{liang2023gsir,shi2023gir,gao2023relightable} and dynamic 3D scene reconstruction \cite{wu20234d,yang2023real}. 



\subsection{Relighting and Inverse Rendering}
Inverse rendering \cite{sato1997object,sato2003illumination} aims to decompose the image’s appearance into the geometry, material properties, and lighting conditions. Most traditional methods simplified the problem by assuming controllable lighting conditions \cite{yu2019inverserendernet,azinovic2019inversepathtracing}. Works based on Nerf explore more complex lighting models to cope with realistic scenarios and extensively utilize MLPs to encode lighting and materials properties. NeRV \cite{srinivasan2021nerv} and Invrender \cite{invrender} train an additional MLP to model the light visibility. NeILF \cite{yao2022neilf} expresses the incident lights as a neural incident light field. NeILF++ \cite{zhang2023neilf++}  integrates VolSDF with NeILF and unifies incident light and outgoing radiance. TensoIR \cite{jin2023tensoir} introduces 
TensoRF representation which enables the computation of visibility and indirect lighting by raytracing. Works based on 3DGS \cite{gao2023relightable,shi2023gir,liang2023gsir} have significantly accelerated training and rendering, enabling real-time relighting and editing. However, these methods still face challenges in real-time high-quality indirect lighting computation, dynamic relighting, and shadow estimation. 
\subsection{Precompute Radiance Transfer}
The fundamental concept of Precompute Radiance Transfer (PRT) \cite{sloan2002precomputed} involves selecting an angular basis comprised of continuous functions, notably Spherical Harmonics (SH), and conducting all light transport operations within this domain. However, Spherical Harmonics are limited in terms of high frequencies, Sloan \cite{liu2004all} then replace it with Haar wavelet. Kristensen \cite{kristensen2005precomputed} further extends PRT to local lighting and pan \cite{pan2007precomputed} extends it to dynamic scenes. Recently there has been a growing interest in using deep learning tools within traditional PRT frameworks \cite{li2019deep,rainer2022neural}, and ideas from PRT have been used in the context of the radiance field \cite{lagunas2021single}. 
